\section{Blockchain Computations}
This section details the decentralized consensus calculations, block validation processes, and cryptographic verification mechanisms. It outlines how traditional blockchain concepts are adapted to serve the simulation's requirements, focusing on the logic and mathematics behind the computations.

\subsection{Nodes and Transactions (Interventions)}
In standard blockchains, nodes are arbitrary computers, and transactions are financial transfers. In this simulation, the network nodes are the countries themselves, meaning the size and power of the network strictly correspond to the geopolitical entities active in the simulation.

The equivalent of blockchain transactions are exogenous system changes called ``God Interventions.'' These interventions dictate forced changes to the global state before the game-theoretic solver attempts to find an equilibrium. The transaction types include:
\begin{itemize}
    \item \texttt{GOD\_INTERVENTION}: Direct adjustments to a country's troop, gold, or population counts.
    \item \texttt{COUNTRY\_ADD} / \texttt{COUNTRY\_REMOVE}: Dynamically provisioning or terminating a node (country) within the network.
    \item \texttt{BUILD\_CASTLE} / \texttt{DEMOLISH\_CASTLE}: Modifying physical fortifications (Tiers 1--3).
    \item \texttt{SET\_TAX\_RATE}: Adjusting the economic tax rate, which indirectly impacts public happiness and income.
\end{itemize}
These transactions are grouped into a mempool and broadcasted to all active country nodes to initiate the mining phase.

\subsection{Block Structure and Stored Data}
A mined block in this architecture stores both the causal transactions and the resultant deterministic state. Specifically, the block payload ($P_{\text{payload}}$) encompasses:
\begin{itemize}
    \item The raw interventions processed from the mempool.
    \item The updated economic ledgers (troops, gold, population, castles, tax, happiness, and rivals).
    \item The calculated geopolitical equilibrium, including predicted alliances and the overarching alliance stability score.
    \item Crisis outcomes such as economic deaths and unhappy emigration.
\end{itemize}
This exhaustive data structure ensures that any node joining the network can verify both the initial state changes (the interventions) and the validity of the solver's complex game-theoretic output.

\subsection{Merkle Root Construction}
To secure the integrity of the data within the block, the system computes a Merkle Root ($M$). In traditional blockchains, a Merkle tree is used to hash pairs of transactions. In this project, to enforce the deterministic state machine, the Merkle Root acts as a unified state commitment hash.

The calculation serializes both the queued mempool interventions and the entirely resolved block payload ($P_{\text{payload}}$) into a JSON string, ensuring deterministic ordering. It then applies a double SHA-256 hash function to this concatenated data:
\begin{equation}
M = \text{SHA-256}\left(\text{SHA-256}\left( \text{JSON}\left( \text{Mempool} \cup P_{\text{payload}} \right) \right)\right)
\end{equation}
This mathematical commitment guarantees that if any node attempts to alter a transaction or manipulate the game-theoretic alliance calculations, the resulting Merkle Root $M$ will drastically change, causing the network to instantly reject the block.

\subsection{Cryptographic Hashing Logic}
The core of the cryptographic verification is the block header, which summarizes the entire state of the block into a single string. The header is constructed by concatenating several critical variables: a version integer, the previous block hash ($H_{\text{prev}}$), the computed Merkle Root ($M$), a UNIX timestamp ($t$), the global difficulty parameter ($D$), the cryptographic nonce ($N$), the winning miner's ID ($\text{ID}_{\text{miner}}$), and the base block reward ($R$):
\begin{equation}
\text{Header} = \text{Concat}\left(1, H_{\text{prev}}, M, t, D, N, \text{ID}_{\text{miner}}, R \right)
\end{equation}
The final cryptographic hash of the block is generated by applying the SHA-256 algorithm twice to the header string:
\begin{equation}
\text{Hash} = \text{SHA-256}\left(\text{SHA-256}\left( \text{Header} \right)\right)
\end{equation}
This resulting 256-bit hexadecimal string serves as the unique identifier for the block and is the value evaluated during the Proof-of-Work phase.

\subsection{Asymmetric Proof-of-Work (PoW)}
Proof-of-Work is a consensus mechanism where nodes expend computational effort to iteratively guess a nonce until they find a hash that meets a specific mathematical condition. This secures the network against spam and malicious attacks by making block generation computationally expensive.

This simulation implements a specialized Asymmetric Proof-of-Work. In traditional blockchains, a node's probability of mining a block is dictated by its hardware computational power, commonly known as its ``hash rate.'' In this architecture, a node's hash rate is mathematically substituted by its geopolitical military strength (i.e., its raw troop count $T_i$). Thus, the mining power $P_i$ serves as the direct equivalent of the node's hash rate. 

A candidate block is considered valid only if its integer hash value is less than or equal to a target threshold $\tau_i$. The target threshold is dynamically calculated for each country:
\begin{equation}
\tau_i = \min\left( T_{\max}, \left( \lfloor \frac{T_{\max}}{D} \rfloor \right) \cdot \max(P_i, 1) \right)
\end{equation}
where $T_{\max} = 2^{256} - 1$ represents the absolute maximum possible hash value. 

Because the target threshold $\tau_i$ is linearly proportional to the mining power $P_i$, countries with massive armies are granted a significantly higher (easier) target threshold. This mathematical bias ensures that military superpowers solve the Proof-of-Work puzzle much faster than weaker nations, directly linking in-game military dominance to blockchain network dominance.

\subsection{Gossip Protocol and Consensus Finality}
In a decentralized network, nodes must propagate information and resolve discrepancies (forks) when multiple blocks are mined simultaneously. Traditional blockchains utilize Peer-to-Peer (P2P) network propagation and the Longest Chain Rule (Nakamoto Consensus) to achieve finality. This simulation adapts these mechanisms for a concurrent threaded environment.

\textbf{Gossip Protocol and Peer Validation:} The active mining country nodes ($\mathcal{A}_{\text{miners}}$) operate concurrently as isolated execution threads. To simulate network propagation, the system implements a shared, thread-safe memory cache known as the gossip cache. As soon as a node $i$ computes a valid hash satisfying $\tau_i$, it writes the candidate block to this shared cache. Concurrently, all other mining threads periodically halt their nonce iteration to inspect the cache. 

If a competing node discovers a gossiped block, it does not blindly trust it. Instead, it instantly subjects the block to a rigorous mathematical verification sequence: it independently verifies the previous hash ($H_{\text{prev}}$), recalculates the Merkle Root ($M$), and recomputes the double SHA-256 hash using the peer's nonce $N$ to ensure it satisfies the peer's specific target threshold. If verification succeeds, the competing node abandons its own PoW loop and "yields" to the verified block, effectively terminating redundant network computation.

\textbf{Consensus Finality (Fork Resolution):} Rather than allowing temporary forks to resolve over multiple blocks via the Longest Chain Rule, the simulation employs a strict First-to-Submit Wins (FTSW) consensus model enforced by the central API gateway. The first verified block submitted to the orchestrator instantly finalizes the transition step. Any subsequent submissions for the same block index are rejected. This guarantees absolute state finality within a single block round, ensuring the simulation UI updates instantly without rollbacks.
